\documentclass[11pt]{article}
\usepackage[margin=1in]{geometry}
\usepackage{amsmath,amssymb,amsthm}
\usepackage{graphicx}
\usepackage[hidelinks]{hyperref}
\setlength{\emergencystretch}{2em}

\newtheorem{theorem}{Theorem}
\newtheorem{lemma}{Lemma}
\newtheorem{corollary}{Corollary}

\newcommand{\B}{\mathcal B}
\newcommand{\K}{\mathcal K}
\newcommand{\C}{\mathbb C}

\title{Stabilized Unital Coronas Fail the Ball-Covering Property}
\author{Run \texttt{fa\_banach\_001}, lane 19}
\date{June 17, 2026}

\begin{document}
\maketitle

\section*{Source Problem}

The source paper is Jinghao Huang, Karimbergenov Kudaybergenov, and Rui Liu,
\emph{The ball-covering property of von Neumann algebras and noncommutative
symmetric spaces}, arXiv:2311.14257.  After proving that a commutative unital
\(C^*\)-algebra \(C(K)\) has BCP/UBCP exactly when \(K\) has a countable
\(\pi\)-basis, the authors ask for a noncommutative version:

\begin{center}
\includegraphics[width=0.92\textwidth]{figures/open_problem_crop.png}
\end{center}

\section*{Statement}

Let \(H\) be an infinite-dimensional separable Hilbert space.  Write
\(\B(H)\) for the bounded operators and \(\K(H)\) for the compact operators.
All tensor products below are minimal/spatial tensor products.

\begin{theorem}
\label{thm:stabilized-corona}
Let \(D\) be a nonzero unital \(C^*\)-algebra.  Then the \(C^*\)-algebra
\[
        \bigl(\B(H)\otimes D\bigr)\big/\bigl(\K(H)\otimes D\bigr)
\]
fails BCP.
\end{theorem}

\begin{corollary}
For every nonzero unital \(C^*\)-algebra \(D\), the multiplier corona of the
stable algebra \(\K(H)\otimes D\) fails BCP.
\end{corollary}

\begin{proof}
For unital \(D\), the standard multiplier identification gives
\[
        M(\K(H)\otimes D)\cong \B(H)\otimes D,
\]
so the multiplier corona is exactly the quotient in
Theorem~\ref{thm:stabilized-corona}.
\end{proof}

\section*{Idea of the Proof}

The Calkin quotient norm is an essential norm: it is seen on right tails after
removing finite-dimensional pieces of \(H\).  The same is true with unital
coefficients \(D\).  Given a countable family of proposed centers, recursively
choose finite-dimensional domain projections \(q_k\) on which the assigned
center nearly realizes its quotient norm.

Since \(Tq_k\) is a finite-column compact-coefficient operator, its range can
be almost captured by some finite-dimensional projection.  Choose a fresh
finite-dimensional range projection \(p_k\) orthogonal to that captured range,
and a partial isometry \(v_k:q_kH\to p_kH\).  Then \(v_k\otimes1_D\) is almost
as far from \(Tq_k\) as \(Tq_k\) is large.  The strong sum \(V=\sum_k v_k\) is
a norm-one partial isometry.  Its corona class has norm one and escapes every
ball \(B(\xi,\|\xi\|)\) in the proposed countable cover.

\section*{Proof}

Let
\[
        \mathcal A=\B(H)\otimes D,\qquad
        \mathcal J=\K(H)\otimes D,
\]
and let \(Q:\mathcal A\to \mathcal A/\mathcal J\) be the quotient map.  If
\(e\) is a finite-rank projection on \(H\), we write \(e\) also for
\(e\otimes 1_D\).

\begin{lemma}[quotient norm on right tails]
\label{lem:quotient-tail}
If \(P_n\) is any increasing sequence of finite-rank projections on \(H\) with
\(P_n\to I\) strongly, then for every \(T\in\mathcal A\),
\[
        \|Q(T)\|=\lim_{n\to\infty}\|T(1-P_n)\|.
\]
\end{lemma}

\begin{proof}
Since \(TP_n\in\mathcal J\), we have
\[
        \|Q(T)\|\le \|T(1-P_n)\|
\]
for every \(n\).  Conversely, if \(C\in\mathcal J\), then
\(\|C(1-P_n)\|\to0\), and hence
\[
        \limsup_n\|T(1-P_n)\|
        \le
        \limsup_n\|(T+C)(1-P_n)\|+\limsup_n\|C(1-P_n)\|
        \le \|T+C\|.
\]
Taking the infimum over \(C\in\mathcal J\) gives the reverse inequality.
\end{proof}

\begin{lemma}[finite right-tail recovery]
\label{lem:finite-tail}
Let \(T\in\mathcal A\), let \(r=\|Q(T)\|\), let \(f\) be a finite-rank
projection on \(H\), and let \(\varepsilon>0\).  There is a nonzero
finite-rank projection \(q\le1-f\) such that
\[
        \|Tq\|>r-\varepsilon.
\]
\end{lemma}

\begin{proof}
Since \(Tf\in\mathcal J\), \(r\le \|T(1-f)\|\).  Represent \(D\)
faithfully on a Hilbert space \(K_D\), so \(\mathcal A\) acts on
\(H\otimes K_D\).  Choose a unit vector
\(\xi\in (1-f)H\otimes K_D\) with
\[
        \|T(1-f)\xi\|>r-\varepsilon/2.
\]
Vectors with finite \(H\)-support are dense in \((1-f)H\otimes K_D\), so
\(\xi\) may be chosen with finite \(H\)-support after a perturbation small
enough to preserve the weaker lower bound \(r-\varepsilon\).
Let \(q\le1-f\) be the projection onto that finite-dimensional \(H\)-support.
Then \(q\xi=\xi\), and therefore \(\|Tq\|>r-\varepsilon\).
\end{proof}

\begin{lemma}[fresh range escape]
\label{lem:fresh-range}
Let \(q\) be a nonzero finite-rank projection on \(H\), let
\(X\in\mathcal A q\), let \(g\) be a finite-rank projection on \(H\), and let
\(\varepsilon>0\).  There are a finite-rank projection \(p\perp g\) and a
partial isometry \(v\in\B(H)\) such that
\[
        v^*v=q,\qquad vv^*=p,
\]
and
\[
        \|(v\otimes1_D)-X\|\ge \|X\|-\varepsilon.
\]
\end{lemma}

\begin{proof}
Since \(q\) is finite rank and \(X=Xq\), we have \(X\in\K(H)\otimes D\).
Choose a
finite-rank projection \(s\) on \(H\) such that
\[
        \|(1-s)X\|<\varepsilon.
\]
Because \(H\) is infinite dimensional, choose a finite-rank projection
\(p\) orthogonal to the span of the ranges of \(g\) and \(s\), with
\(\operatorname{rank}p=\operatorname{rank}q\), and
choose a partial isometry \(v\in\B(H)\) with \(v^*v=q\) and \(vv^*=p\).
Then \(pX=p(1-s)X\), so \(\|pX\|<\varepsilon\).  Since
\((1-p)(v\otimes1_D)=0\),
\[
        \|(v\otimes1_D)-X\|
        \ge
        \|(1-p)((v\otimes1_D)-X)\|
        =
        \|(1-p)X\|
        \ge
        \|X\|-\|pX\|
        >
        \|X\|-\varepsilon.
\]
\end{proof}

\begin{proof}[Proof of Theorem~\ref{thm:stabilized-corona}]
Let \(\xi_i=Q(T_i)\), \(i\ge1\), be an arbitrary countable family of centers
in \(\mathcal A/\mathcal J\), with \(T_i\in\mathcal A\).  Put
\[
        r_i=\|\xi_i\|.
\]
Choose a map \(\pi:\mathbb N\to\mathbb N\) such that every value occurs
infinitely often.

We recursively choose finite-rank domain projections \(q_k\), finite-rank
range projections \(p_k\), and partial isometries \(v_k\in\B(H)\).  Suppose the
previous choices have been made, and let \(g_{k-1}\) be a finite-rank
projection dominating all previous \(q_j\)'s and \(p_j\)'s.  Put
\(i=\pi(k)\).  By Lemma~\ref{lem:finite-tail}, choose a nonzero finite-rank
projection \(q_k\le1-g_{k-1}\) such that
\[
        \|T_iq_k\|>r_i-\frac1k.
\]
Apply Lemma~\ref{lem:fresh-range} to \(q=q_k\), \(X=T_iq_k\), and
\(g=g_{k-1}\vee q_k\), obtaining a projection \(p_k\perp(g_{k-1}\vee q_k)\)
and a partial isometry \(v_k\) with
\[
        v_k^*v_k=q_k,\qquad v_kv_k^*=p_k,
\]
such that
\[
        \|(v_k\otimes1_D)-T_iq_k\|>r_i-\frac2k.
\]
This makes the \(q_k\)'s pairwise orthogonal, the \(p_k\)'s pairwise
orthogonal, and every \(p_k\) orthogonal to every \(q_j\).

Let \(V\in\B(H)\) be the strong sum of the partial isometries \(v_k\).  Since
the initial projections \(q_k\) are pairwise orthogonal and the final
projections \(p_k\) are pairwise orthogonal, \(V\) is a partial isometry of
norm one.  Set
\[
        B=V\otimes1_D\in\mathcal A.
\]
The projections \(q_k\) converge strongly to \(0\), so for every
\(C\in\mathcal J\),
\[
        \|Cq_k\|\longrightarrow0.
\]
Indeed, this is immediate for elementary tensors \(C_0\otimes d\) with
\(C_0\in\K(H)\), and then follows for all \(C\in\K(H)\otimes D\) by norm
approximation.
Therefore \(\|Q(B)\|=1\), because \(\|B\|\le1\) and
\[
        \|B-C\|\ge \|(B-C)q_k\|\ge \|Bq_k\|-\|Cq_k\|
        =1-\|Cq_k\|
\]
for all \(k\).

Fix \(i\).  If \(r_i=0\), then \(\|Q(B)-\xi_i\|\ge0=r_i\).  If \(r_i>0\),
take \(k\to\infty\) along the infinite set \(\{k:\pi(k)=i\}\).  For any
\(C\in\mathcal J\),
\[
        \|B-T_i-C\|
        \ge
        \|(B-T_i-C)q_k\|
        \ge
        \|Bq_k-T_iq_k\|-\|Cq_k\|
        >
        r_i-\frac2k-\|Cq_k\|.
\]
Letting \(k\to\infty\) along this subsequence gives
\[
        \|B-T_i-C\|\ge r_i.
\]
Taking the infimum over \(C\in\mathcal J\) yields
\[
        \|Q(B)-\xi_i\|\ge r_i=\|\xi_i\|.
\]
Thus \(Q(B)\) is a unit vector outside every open ball
\[
        B(\xi_i,\|\xi_i\|).
\]
Since the countable family of centers was arbitrary, the quotient fails BCP.
\end{proof}

\section*{Scope and Relation to Earlier Packets}

For \(D=\C\), Theorem~\ref{thm:stabilized-corona} recovers the Calkin-algebra
obstruction.  The point of the present packet is that the same proof works
with arbitrary nonzero unital coefficient algebra \(D\).  This is different
from the product and reduced-product packets: the quotient is a genuine
operator-algebra corona, and the proof uses finite-rank right-tail recovery
instead of coordinatewise reduced-product norms.

The theorem does not characterize arbitrary \(C^*\)-algebras with BCP.  It
places a broad stabilized corona family on the negative side and supplies a
concrete model for future multiplier-corona attacks.

\section*{Verification and Search Bounds}

The run indexes were searched for arXiv:2311.14257 together with stabilized
corona, \(B(H)\otimes D\), \(K(H)\otimes D\), Calkin, multiplier, and
partial-isometry keywords.  No existing packet for this coefficient-stabilized
corona obstruction was found.  The proof is self-contained modulo standard
facts about compact-operator tensor ideals and the multiplier identification
for \(\K(H)\otimes D\) when \(D\) is unital.

\section*{Human Review Focus}

Check Lemma~\ref{lem:finite-tail}, especially the finite \(H\)-support
approximation after representing \(D\) faithfully, and Lemma~\ref{lem:fresh-range},
where the fresh range projection is chosen away from the finite-column image.
The final diagonal argument then follows the same shape as the previous
product-quotient obstruction.

\begin{thebibliography}{9}

\bibitem{HKL}
Jinghao Huang, Karimbergenov Kudaybergenov, and Rui Liu,
\emph{The ball-covering property of von Neumann algebras and noncommutative
symmetric spaces},
arXiv:2311.14257, 2023.

\end{thebibliography}

\end{document}
